\documentclass[a4paper, oneside]{article}
\usepackage{graphicx}
\usepackage{amsthm}
\usepackage{amsmath}
\usepackage{amssymb}
\usepackage[a4paper,
            bindingoffset=0.2in,
            left=2cm,
            right=2cm,
            top=2cm,
            bottom=2cm,
            footskip=.25in]{geometry}
\usepackage[italian]{babel}
\usepackage{pgfplots}
\usepackage{tabularx}
\usepackage{tikz}
\usepackage{wrapfig}
\usepackage{color}
\usepackage[d]{esvect}
\definecolor{page}{rgb}{0.129,0.157,0.212}
\pagecolor{page}
\color{white}
\graphicspath{ {./images/} }
\usetikzlibrary{shapes.geometric}
\usetikzlibrary{datavisualization}
\usetikzlibrary{datavisualization.formats.functions}
\usetikzlibrary{patterns}
\pgfplotsset{width=10cm,compat=1.18}

\title{Appunti di Metodi}
\author{Tommaso Miliani}
\date{04-02-26}

\begin{document}
\newtheoremstyle{theoremEnv}
                {}          % Space above
                {}          % Space below
                {\slshape}  % Body font
                {}          % Indent amount
                {\bfseries} % Head font
                {.}         % Punctuation after head
                {\newline}  % Space after theorem head
                {}          % Theorem head spec
\theoremstyle{theoremEnv}

\newtheorem{definition}{Definizione}[section]
\newtheorem{theorem}{Teorema}[section]
\newtheorem{lemma}{Proposizione}[section]
\newtheorem{observation}{Osservazione}[section]
\newtheorem{corollary}{Corollario}[theorem]
\newtheorem{example}{Esempio}[section]
\newtheorem{remark}{Enunciato}[section]

\maketitle

\section{Operazioni elementari sulle serie di potenze}
La prima è la \textbf{somma di due serie di potenze}: ossia la somma termine a termine
\begin{gather*}
    f_a = \sum a_n z^{n} \qquad f_b = \sum b_nz^{n}
\end{gather*}
Allora 
\begin{gather*}
    (f_a + f_b)(z) = f_a(z) + f_b(z) \ \Longrightarrow \ \sum (a_n + b_n)z^{n}
\end{gather*}
Questa uguaglianza vale nell'intersezione dei domini di convergenza. Preso $z \in C_1 \cap C_2$, 
allora questa uguaglianza vale se $C_1$ è il dominio di convergenza di $f_a$ e $C_2$ quello 
di $f_b$. Questo dominio è $\geq$ del dominio di convergenza più piccolo.
Si definisce anche il prodotto per uno scalare:
\begin{gather*}
    \lambda f(z) = \lambda f(z)
\end{gather*}
Dunque
\begin{gather*}
    \sum \lambda a_n z^{n}
\end{gather*}
Il raggio di convergenza di questa serie è lo stesso a meno che $\lambda = 0$, altrimenti 
il raggio di convergenza è infinito. Si può anche definire il prodotto tra due serie:
\begin{gather*}
    (f_a \cdot f_b) (z) = f_a(z) \cdot f_b(z) 
\end{gather*}
Con $z \in C_1 \cap C_2$. I coefficienti di questa serie prodotto non sono 
semplicemente il prodotto dei coefficienti. I nuovi coefficienti $c_n$ sono 
\begin{gather*}
    c_n = \sum_{k = 0}^{n}a_kb_{n - k} 
\end{gather*}
Nella serie originale infatti si ha che
\begin{gather*}
    a_k z^{k} b_{n - k} b_{n - k} z^{n - k} = a_k b_{n - k}z^{n}
\end{gather*}
Essenzialmente si riordinano i termini della serie.  Si può anche vedere
come
\begin{gather*}
    \sum_{l = 0}^{\infty } a_l z^{l} \left(\sum_{m = 0}^{} \right) 
\end{gather*}
È come si si organizzasse le serie in una matrice e si ottenesse la traccia
delle sottomatrici. 
\begin{example} [Esempio di applicazione del prodotto]
    Si vuole dimostrare
    \begin{gather*}
        \sin(2x) = 2\sin x \cos x
    \end{gather*}
    utilizzando la formula di prodotto e lo sviluppo in serie. Si 
    sviluppa il secondo termine:
    \begin{gather*}
        2\left(\sum_{n = 1}^{\infty } \frac{(-1)^{n} x^{2n + 1}}{(2n + 1)!} \right) \left(\sum_{m = 1}^{\infty } \frac{(-1)^{m}x^{2m}}{2m!} \right)
    \end{gather*}
    Dato che compaiono potenze dispari e pari, quando si fanno i prodotti verranno solamente 
    delle potenze dispari, dunque, chiamato $C_{2k + 1}$ come i termini dispari:
    \begin{gather*}
        C_{2k + 1} = 2\sum_{n + m = k}^{} \frac{(-1)^{n}}{(2n + 1) !} \frac{(-1)^{m}}{2m!} 
    \end{gather*}
    Si moltiplica e divide per dei fattoriali:
    \begin{gather*}
        = \frac{2(-1)^{k}}{(2k + 1) !} \sum_{m + n = k} \frac{(2k + 1)!}{(2n + 1)!(2m)!} \ \Longrightarrow \ \frac{(2k + 1)!}{(2n + 1)! (2m)!} = \begin{pmatrix} 2k + 1 \\
        2n + 1 \end{pmatrix} =  \begin{pmatrix} 2k + 1  \\
        2m\end{pmatrix} 
    \end{gather*}
    Ottenendo lo sviluppo del binomio, si può ora semplificare
    \begin{gather*}
        \frac{(-1)^{k}}{(2k + 1)!} \left(\sum_{n = 0}^{k} \begin{pmatrix} 2k + 1 \\
        2n + 1 \end{pmatrix} + \sum_{m = 0}^{k} \begin{pmatrix} 2k + 1\\
        2m \end{pmatrix}  \right)
    \end{gather*}
    Questa uguaglianza è come se si sommassero i coefficienti del binomio su tutti gli indici da
    zero a $2k + 1$ dato che si sta sommando sia dagli indici pari che dagli indici dispari
    \begin{gather*}
        \sum_{i = 0}^{2k + 1} \begin{pmatrix} 2k + 1 \\
        i \end{pmatrix}  = 2^{2k + 1}
    \end{gather*}
    L'ultima uguaglianza si può dimostrare semplicemente visualizzandola come l'espansione del seguente:
    \begin{gather*}
        (1 + 1)^{2k + 1}
    \end{gather*}
    In altre parole 
    \begin{gather*}
        (a + b)^{N} = \sum_{k = 0}^{N} a^{k}b^{N - k} \begin{pmatrix} N \\
        k \end{pmatrix} 
    \end{gather*}
    A questo punto diventa
    \begin{gather*}
        \frac{(-1)^{k} 2^{2k + 1}}{(2k + 1)!} = C_{2k + 1}
    \end{gather*}
    Ossia l'espansione in serie del prodotto tra seno e conseno.
\end{example}

\subsection{Derivazione delle serie}
Ricordando che siamo su $\mathbb{C}$, si deve dunque estendere la definizone di 
derivata al caso di funzione complessa. La derivata complessa è più forte della derivata a
reale.
\begin{definition}[Derivazione]
    Supponendo di avere una funzione $f : \mathbb{D} \to \mathbb{C}$ tale per cui $z \in \mathbb{D}$ sia punto di 
    accumulazione. La derivata complessa di $f$ è dunque definita come 
    \begin{align}
        f'(z) \equiv \lim_{w \to z} \frac{f(w) - f(z)}{w - z} \qquad w \in \mathbb{D} 
    \end{align}
    Quando si è su $\mathbb{C}$ questo limite si può fare da qualunque direzione si voglia. L'importante 
    è che questo limite esista per tutte le direzioni prese e che sia uguale per tutte le direzioni. 
\end{definition}
\begin{theorem}[Derivabilità $\ \Longrightarrow \ $ continuità]
    Come per le funzioni su $\mathbb{R}$, se $f$ è derivabile in $z$, allora è continua in $z$. 
\end{theorem} \noindent
Proprietà delle derivate dei polinomi complessi:
\begin{align*}
    &\frac{d}{dx}(z^{n}) = nz^{n - 1} \quad n \geq 1 \\
    &\frac{d}{dz}1 = 0
\end{align*}
Con $n \in \mathbb{Z}$ e $z \in \mathbb{C}$. Si potrà poi estendere anche al caso in cui $n \in \mathbb{C}$.
Dato che si sta facendo la derivata di una serie di potenza, si farà termine a termine 
e ci si chiede quale sarà il raggio di convergenza della serie derivata.
\begin{theorem}[Il raggio di convergenza derivato è uguale a quello di partenza]
    Il raggio di convergenza della serie derivata è uguale al raggio di convergenza della serie 
    iniziale.
\end{theorem}
\begin{proof}
    Utilizzando il criterio del confronto, imponendo che sia $\leq$ e $\geq$. Chiamati 
    $R_a'$ e $R_a$ i raggi di convergenza della serie derivata e di quella iniziale. \\
    $R_a' \leq R_a$. Preso un numero $z \in \mathbb{C}$ : $\left| z \right| < R_a'$. La
    serie delle derivate è 
    \begin{gather*}
        f' = \sum_{n = 1}^{\infty } na_n z^{n - 1} \qquad f = \sum_{n = 0}^{\infty }a_nz^{n}  
    \end{gather*} 
    Si vuole allora maggiorare:
    \begin{gather*}
        \left| a_n z^{n} \right| = \left| z \right| \left| a_n z^{n - 1} \right| \leq n \left| a_n z^{n - 1} \right| \qquad \text{se } n \geq \left| z \right|     
    \end{gather*}
    Per qualunque $z$ si può trovare un $n$ sufficientemente grande tale per cui vale
    l'ultima uguaglianza per tutti gli $n$ che vengono dopo. Dunque se converge questa serie maggiorata
    con i termini della serie derivata, allora convergerà anche quella di partenza e quella derivata e 
    si è dimostrato che $R_a' \leq R_a$.  \\
    $R_a' \geq R_a$. Preso sempre $\left| z \right| < R_a$, si prende allora un punto in cui 
    si sa che la serie originale converge. Esiste dunque 
    \begin{gather*}
        \exists \ w : \left| z \right| < \left| w \right| < R_a  
    \end{gather*} 
    Si può dunque scrivere
    \begin{gather*}
        \left| na_nz^{n - 1} \right| = n\left| a_n \right| \left| w \right|^{n - 1}\left| \frac{z}{w} \right|    ^{n -1}
    \end{gather*}
    Si può dunque ottenere, riorganizzando i termini
    \begin{gather*}
        = \left(\frac{n}{\left| w \right| } \left| \frac{z}{w} \right|^{n - 1} \right)\left| a_n w^{n} \right| 
    \end{gather*}
    Ci si è ricondotti ad una uguaglianza in cui ci sono i termini della serie originaria semplicemente 
    espressa con il termine di convergenza $w$. I termini dentro la parentesi sono 
    una serie geometrica che è minore stretto di uno per come è stato scelto $w$. Dunque 
    quella roba è sicuramente limitata: per $n \to \infty$ tende a zero. Si è dunque 
    dimostrato che i termini della serie derivata sono maggiorati dai termini della serie 
    iniziali da un $w$ dentro al raggio di convergenza. Allora convergono entrambe e 
    si è dimostrato che $R_a \leq R_a'$. 
\end{proof}

\begin{theorem}[Teorema di derivazione termine a termine]
    Data 
    \begin{gather*}
        f(z) = \sum_{n }^{} a_n z^{n}
    \end{gather*}
    Allora 
    \begin{gather*}
        f'(z) = \sum_{n }^{} na_n z^{n - 1}
    \end{gather*}
    Dunque si definisce la derivata della serie di potenze come 
    la somma delle derivate dei singoli elementi che è equivalente a fare la derivata della funzione somma.
    Così come per le funzioni reali, anche per queste funzioni 
    \begin{gather*}
        f(0) = 0 \qquad f'(0) = a_1 \qquad f''(0) = 2a_2 \qquad f'''(0) = 3! a_3
    \end{gather*}
    Dunque in generale
    \begin{gather*}
        f^{(n)}(0) = n!a_n \ \Longrightarrow \ a_n = \frac{1}{n!}f^{(n)}( 0)
    \end{gather*}
    Che non è altro che l'espansione di Taylor in serie
    \begin{gather*}
        f(z) = \sum_{n = 0}^{\infty } \frac{f^{(n)}(0)}{n!}z^{n}
    \end{gather*}
    Ossia la generalizzazione su $\mathbb{C}$ di quello che si è già visto su 
    $\mathbb{R}$. 
\end{theorem}\noindent
Supponendo di avere $f(z) = \sum_{z }^{} a_n z^{n}$ e $f(z) = \sum_{z} b_n z^{n}$ 
su di un certo dominio non banale $C_a$. Si può concludere che se $R_a > 0$ allora $a_n = b_n$. 
Basta dunque conoscere $f$ in un cerchio arbitrariamente piccolo intorno all'origine, che basta per determinare 
tutte le derivate vicino allo zero, dunque si conosce $a_n$. Questo non vale se $R_a = 0 $. 
Se ci si avvicina al bordo radialmente, la funzione è continua ma non è detto che lo sia 
se si procede tangenzialmente al bordo (teorema di Abbel).

\begin{theorem}[Teorema di integrazione termine a termine]
    Avendo la solita serie di potenze con raggio di convergenza $0 < R_a \leq + \infty$:
    \begin{gather*}
        \sum_{n }^{} a_n z^{n}
    \end{gather*}
    Si vuole adesso definire, integrando termine a termine, 
    \begin{gather*}
        c_0 + \sum_{n = 0}^{\infty } \frac{a_n}{n + 1}z^{n + 1}  
    \end{gather*}
    L'integrale è dunque definito a meno di una costante quando 
    si integra termine a termine. Per questo si è riassunto $c_0$ come 
    il punto di inizio di integrazione che riassume tutte le costanti.
    Il raggio di convergenza è lo stesso di quella di partenza poiché, derivando,
    si torna al teorema di prima. La derivata in senso complesso è la funzione iniziale e 
    hanno anche una primitiva.
\end{theorem}
\begin{example}[Serie logaritmo]
\begin{gather*}
    \sum (-1)^{n}z^{n} = \frac{1}{1 + z}  \\
    c + \sum_{n = 0}^{\infty } (-1)^{z} \frac{z^{n + 1}}{n + 1} = c + \sum_{m = 1}^{\infty } (-1)^{m + 1} \frac{z^{m}}{m} = c + \log(1 + z) 
\end{gather*}
\end{example}


\section{Definire funzioni con le serie di espansione}
Tramite le serie si vuole definire una delle funzioni più importanti 
per l'analisi complessa, armonica e spazi di Hilbert ed è la 
\textbf{funzione esponenziale}. Questa funzione ha la peculiarità che
la sua derivata ed il suo integrale sono la funzione stessa. questa definisce 
anche il seno ed il coseno, così come il logaritmo (inversa dell'esponenziale).
\begin{definition}[Funzione esponenziale]
    Si definisce in maniera formale la funzione epsonenziale utilizzando
    le equazioni differenziali:
    \begin{gather*}
        \left\{\begin{array}{l}
            f' = f \\
            f(0) = 1
        \end{array}\right.
    \end{gather*}
    La soluzione unica di questa equazione differenziale è $f = \exp(x)$. 
    Si può scrivere anche sui complessi con la derivata complessa e
    dunque la soluzione sui complessi sarà $\exp (z)$, $z \in \mathbb{C}$.
\end{definition}
A partire da questa definizione non si riesce a calcolare (con un
certo grado di approssimazione) la funzione esponenziale, dovrei utilizzare
delle serie
\begin{definition}[Funzione esponenziale per serie]
    Scrivendo in forma generica una serie
    \begin{gather*}
        \exp(z) = \sum_{n}^{} a_nz^{n} = 1 + a_1 z+ \dots
    \end{gather*}
    Si può dunque utilizzare l'espansione in serie per risolvere le 
    equazioni differenziali:
    \begin{gather*}
        \exp'(z) = \sum_{n }^{} na_nz^{n - 1} = a_1 + 2a_2 z + 3a_3 z^{2} + \dots + na_nz^{n - 1}
    \end{gather*}
    uguagliando le due, si può ottenere
    \begin{gather*}
        a_1 = 1 \qquad a_2 = \frac{a_1}{2} \qquad a_3 = \frac{a_2}{3} \quad\dots\quad a_n = \frac{1}{n!}
    \end{gather*}
    Quindi la serie che si è trovato è
    \begin{gather*}
        \exp(z) = \sum_{n = 0}^{\infty } \frac{1}{n!}z^{n}
    \end{gather*}
    Questa serie ha raggio di convergenza $R = \infty$, dunque converge su tutto $\mathbb{C}$ 
    ed è una buona definizione dell'esponenziale. Un'altra notazione che si utilizza è
    \begin{gather*}
        \exp(z) = e^{z}
    \end{gather*}
\end{definition} \noindent
Su $\mathbb{C}$ valgono le stesse proprietà che valgono su $\mathbb{R}$:
\begin{enumerate}
    \item $e^{w}\cdot e^{z} = e^{w + z}$, $\forall w, z \in \mathbb{C}$.
    \item $e^{-z} = (e^{z})^{-1}$ ; 
    \item $e^{z^{\star}} = (e^{z})^{\star}$;
    \item $\left| e^{z} \right| = \left| e^{x + iy} \right| = \left| e^{z} e^{z^{\star}} \right|^{\frac{1}{2}} = e^{x}$. 
    \item $\left| e^{iy} \right| = 1$.    
\end{enumerate}
Si potrebbe dimostrare la prima con il prodotto in serie, oppure attraverso
le equazioni differenziali:
\begin{gather*}
    \frac{d}{dz}(e^{w} \cdot e^{z}) = e^{w}\cdot e^{z} = \frac{d}{dz}(e^{w + z})
\end{gather*}
Quando $z = 0$ sono $e^{w} = e^{w}$. Un'altra conseguenza 
diretta di questa formula è la seconda:
\begin{gather*}
    e^{-z}e^{z} = e^{0} = 1
\end{gather*}
Così come sui reali, anche nei complessi non si azzera mai e dunque 
\begin{gather*}
    Im\left(e^{z}\right) = \mathbb{C} - \{0\}
\end{gather*}

\section{Definizione di funzioni trigonometriche tramite l'esponenziale}
Dato $e^{z}$, se $z = x + iy$, se $x = 0$, allora 
si ottiene che il modulo di $e^{z}$ è $\left| e^{iy} \right|$.
Si vuole ora trovare la derivata di questa funzione per vedere
come varia:
\begin{gather*}
    \frac{d}{dy}e^{iy} = ie^{iy} \ \Longrightarrow \ \left| \frac{de^{iy}}{dy} \right| = 0
\end{gather*} 
Al variare di $y$ $e^{iy}$ si sposta di velocità costante $1$ sul cerchio unitario 
nel verso antiorario poiché se $y = 0$, la direzione è $i$. Dato che 
non si annulla mai, al variare di $y$ ci si avvolge infinite volte intorno al cerchio 
percorrendo un arco pari a $y$ in radianti. Dunque
\begin{gather*}
    e^{iy} = e^{i(y + 2\pi k)} \quad k \in \mathbb{Z}
\end{gather*}
Se si volesse una definizione più precisa di $\pi$, si può definire come 
\begin{gather*}
    e^{i\pi} = - 1
\end{gather*}
Ossia $\pi$ è il valore minimo di $y > 0$ tale che vale quella sopra. Questa si dimostra
\begin{gather*}
    e^{2\pi i} = e^{i \pi} e^{i\pi } = (-1)^{2} = 1
\end{gather*}
Dunque aggiungendo $2\pi k $ all'esponenziale, esso rimane uguale. Si è dunque 
scoperta la struttura dell'esponenziale sul piano complesso: esso è periodico con 
periodo $2\pi i$. Qualunque $z$ si prenda, l'esponenziale è periodico su fasce di 
ampiezza $i$. In geometria si definsice seno e coseno come la proiezione lungo $x$
e la proiezione lungo $y$ di un punto su di una circonferenza. 
\begin{definition}[Seno e coseno]
    SI definisce seno e coseno come la proiezione della parte immaginaria e reale dell'esponenziale
    \begin{align}
        e^{iy} = \cos y + i\sin y
    \end{align}
    Se si considera $e^{iy}$ e $e^{iy}$, si può ricavare il seno ed il 
    coseno come combinazione di esponenziali
    \begin{align}
        \cos y = \frac{e^{iy} + e^{-iy}}{2}
    \end{align}
    ed il seno 
    \begin{align}
        \sin y = \frac{e^{iy} - e^{-iy}}{2i}
    \end{align}
    Da queste si vede che il coseno è pari ed il seno dispari. Tramite queste si ricavano 
    le espansioni in serie del seno e del coseno. Quello fatto fino ad ora è per $y$ reale, 
    si può anche definire, dato $z \in \mathbb{C}$, 
    \begin{align}
        \cos z = \frac{e^{iz} + e^{-iz}}{2} 
    \end{align}
    e 
    \begin{align}
        \sin z = \frac{e^{iz} - e^{-iz}}{2i}
    \end{align}
    Sono dunque le estensioni in $\mathbb{C}$ del seno e del coseno e 
    coincidono con le definizioni su $\mathbb{R}$ e a partire da queste 
    si possono anche fare le serie partendo dalla serie esponenziale $e^{iz}$. La 
    serie del coseno contiene solo potenze pari, mentre quelle del seno le 
    potenze dispari. 
\end{definition} \noindent
Partendo dalla serie
\begin{gather*}
    e^{iz} = \sum_{k = 0}^{\infty } (-1)^{k} \frac{z^{2k}}{(2k)!} + i \sum_{k = 0}^{\infty } (-1)^{k} \frac{z^{2k + 1}}{(2k + 1)!}  
\end{gather*}
Si è separato i termini pari ed i termini dispari. Dunque si è ottenuto gli sviluppi per
il coseno ed il seno portando $i^{2k}$ e $i^{2k + 1}$ rispettivamente davanti a tutte e due
(per questo compare $i$ solo per il seno). La serie per il coseno sarà
\begin{gather*}
    \cos z = \sum_{k = 0}^{\infty } (-1)^{k} \frac{z^{2k}}{(2k)!}  
\end{gather*}
Mentre quella per il seno 
\begin{gather*}
    \sin z = \sum_{k = 0}^{\infty } (-1)^{k} \frac{z^{2k + 1}}{(2k + 1)!} 
\end{gather*}
$ i$ scompare perché si divide per $2i$. 
\begin{definition}[Derivate di seno e coseno]
    Attraverso le regole di derivazione per serie si possono anche 
    dimsotrare le regole di derivazione anche su $\mathbb{C}$: 
    \begin{align}
        \frac{d}{dz} \cos z &= - \sin z \\
        \frac{d}{dz} \sin z &= \cos z
    \end{align}
\end{definition} \noindent
Se si calcolasse
\begin{align*}
    \cos iy &= \cosh y \\
    \sin iy &= i \sinh y
\end{align*}
Si può ricavare anche 
\begin{gather*}
    \cos(\alpha \pm \beta ) \\
    \sin(\alpha \pm \beta )
\end{gather*}
E lo stesso per i numeri complessi. Si nota che l'esponenziale complesso,
permette di ottenere i punti sulla circonferenza di raggio 1 in funzione 
dell'angolo $\theta$. Si può estrarre dunque le radici n-esime attraverso l'esponenziale.
 

\end{document}